% !TeX program = pdflatex
% !TeX TXS-program:compile = txs:///pdflatex/[-shell-escape]
% !TeX TXS-program:pdflatex = pdflatex -synctex=1 -interaction=nonstopmode --shell-escape %.tex

\documentclass[../../../main_compilation.tex]{subfiles}

\begin{document}

% ================================================================================
% Encabezado del Laboratorio PortSwigger
% ================================================================================
\section[PortSwigger - TITULO\_LAB]{TITULO\_LAB — PortSwigger Academy}
\label{sec:portswigger-titulo-lab}

\targetSummary%
{TITULO\_LAB}% Nombre del Laboratorio
{https://lab-id.web-security-academy.net}% URL del Laboratorio
{PortSwigger}% Plataforma
{Practitioner}% Dificultad: Apprentice, Practitioner, Expert
{Web Application}% Entorno
{SQLi / XSS / CSRF / SSRF / Auth}% Categoría de Vulnerabilidad
{\today}% Fecha
{Completado}% Estado
{N/A (Web Lab)}% Flag / Solución

\vspace{0.4cm}

% ================================================================================
% 1. Descripción del Laboratorio
% ================================================================================
\subsection{Descripción y Objetivo del Laboratorio}
\textbf{Objetivo:} Describir la meta del laboratorio (por ejemplo, realizar un ataque de SQLi para extraer la contraseña del usuario \texttt{administrator} o ejecutar un XSS almacenado para robar cookies de sesión).

\begin{vulnbox}{Tipo de Vulnerabilidad Web}{High}
	\textbf{Categoría:} OWASP Top 10 — Inyección / Control de Acceso / etc.\\
	\textbf{Mecanismo:} Cómo la aplicación procesa incorrectamente los datos de entrada del usuario.
\end{vulnbox}

% ================================================================================
% 2. Análisis del Tráfico HTTP & Explotación
% ================================================================================
\subsection{Paso a Paso de la Explotación (Burp Suite)}

\begin{terminal}[Burp Suite Repeater / Request]
	POST /login HTTP/1.1
	Host: lab-id.web-security-academy.net
	Content-Type: application/x-www-form-urlencoded

	username=administrator'--&password=cualquiera
\end{terminal}

\begin{notebox}[Payload Utilizado]
	Explicación detallada de la estructura del payload y por qué evade los filtros o resuelve la lógica vulnerable.
\end{notebox}

% ================================================================================
% 3. Remediación Defensiva (Secure Coding)
% ================================================================================
\subsection{Remediación en Código Seguro}

\begin{mitigationbox}[Defensa en Profundidad]
	\begin{itemize}[leftmargin=*]
		\item Uso de consultas parametrizadas / \textit{Prepared Statements}.
		\item Saneamiento y codificación contextual de salidas (\textit{Context-aware output encoding}).
		\item Implementación de cabeceras de seguridad (\textit{CSP}, \textit{SameSite cookies}, \textit{HSTS}).
	\end{itemize}
\end{mitigationbox}

% ================================================================================
% 4. Lecciones y Conceptos Aprendidos
% ================================================================================
\subsection{Conocimientos y Técnicas Aprendidas}

\begin{learningbox}[Lecciones del Laboratorio Web]
	\begin{itemize}[leftmargin=*]
		\item \textbf{Mecanismo de Vulnerabilidad:} Causa raíz del fallo en la lógica o validación del servidor/cliente.
		\item \textbf{Evasión / Filtro Superado:} Técnica empleada para eludir WAF o validaciones de entrada.
	\end{itemize}
\end{learningbox}

\begin{sourcebox}[Fuentes y Referencias Web Security]
	\begin{itemize}[leftmargin=*]
		\resourceItem{Web}{PortSwigger Web Security Academy Materials}{\url{https://portswigger.net/web-security}}
		\resourceItem{Docs}{OWASP Web Security Testing Guide (WSTG)}{\url{https://owasp.org/www-project-web-security-testing-guide/}}
	\end{itemize}
\end{sourcebox}

\end{document}
